% Synthetic diagrams for keys that tests/cd does not reach but crates/commute recognises.  tests/oracle/check_commute_extra.py
% compares commute with pdflatex on them at run time; nothing derived from this file is committed.
\documentclass{article}
\pdfcompresslevel=0 \pdfobjcompresslevel=0
\usepackage{amsmath,amssymb}
\usepackage{tikz-cd}
\usetikzlibrary{decorations.pathmorphing}
\pagestyle{empty}
\begin{document}
\begin{tikzcd} A \arrow[r,densely dotted] & B \arrow[r,loosely dotted] & C \arrow[r,densely dashed] & D \arrow[r,loosely dashed] & E \arrow[r,dash pattern=on 2pt off 1pt on 3pt off 4pt] & F \end{tikzcd}

\begin{tikzcd}[column sep=large] A \arrow[r,"a" {inner xsep=3pt}] & B \arrow[r,"b" {inner ysep=3pt,fill=red}] & C \arrow[r,"c"{outer sep=2pt}] & D \arrow[r,"d"{minimum width=1cm,fill=blue}] & E \arrow[r,"e"{minimum height=1cm,description}] & F \end{tikzcd}

\begin{tikzcd}[column sep=large] A \arrow[r,"a"{rotate=30,xshift=4pt}] & B \arrow[r,"b"{xshift=4pt,rotate=30}] & C \arrow[r,"c"{yshift=2pt,rotate=45,xshift=1pt}] & D \end{tikzcd}

\begin{tikzcd}[sep=large] A \arrow[rd,"c"{sloped,rotate=10}] \arrow[r,"d"{sloped,description}] & B \arrow[d,"e"{sloped,fill=red}] \\ & C \arrow[lu,"f"{allow upside down,sloped}] \end{tikzcd}

\begin{tikzcd}[row sep=large] A \arrow[r,"a"{alias=U}] \arrow[r,bend right=50,"b"'{name=D,alias=W}] & B \arrow[Rightarrow,from=U,to=W] \end{tikzcd}

\begin{tikzcd} A \arrow[r,double=red] & B \arrow[r,transform canvas={shift={(1pt,2pt)}},"x"] & C \arrow[r,"y"',labels={inner sep=2pt}] & D \arrow[r,red,double,"z"] & E \end{tikzcd}

\begin{tikzcd} A \arrow[r,start anchor={[shift={(1pt,2pt)}]north},end anchor={[xshift=1pt,yshift=2pt]south}] & B \end{tikzcd}

\begin{tikzcd} A \arrow[r,bend left,shift left=3pt] & B \arrow[r,out=30,in=150,xshift=3pt] & C \arrow[r,bend right,yshift=-2pt,"f"] & D \arrow[r,start anchor=north,end anchor=north,bend left] & E \end{tikzcd}

\begin{tikzcd} A \arrow[loop] & B \arrow[loop below,"x"] & C \arrow[loop,out=60,in=30,looseness=4] \end{tikzcd}

\begin{tikzcd} A \arrow[r,bend left,out looseness=2] & B \arrow[r,distance=1cm,out=90,in=90] & C \arrow[r,bend right,min distance=8mm] & D \arrow[r,bend left,in looseness=.5,max distance=3pt] & E \end{tikzcd}

{\tikzcdset{arrow style=tikz}\begin{tikzcd} A \arrow[r,harpoon] & B \arrow[r,harpoon'] & C \arrow[r,tail] & D \arrow[r,leftarrow,tail] & E \arrow[r,hook'] & F \arrow[r,thick,two heads] & G \end{tikzcd}}

{\tikzcdset{arrow style=Latin Modern}\begin{tikzcd} A \arrow[r,hook'] & B \arrow[r,Mapsto] & C \arrow[r,hook,thick] & D \arrow[r,Mapsfrom] & E \arrow[d,hookleftarrow] \\ & & & & F \end{tikzcd}}

\begin{tikzcd}[foo/.style={hook,thick},bar/.style={"x"},baz/.style={red},baz/.append style={dashed},qux/.style={near start}] A \arrow[r,foo] & B \arrow[r,bar] & C \arrow[r,baz,"y"{qux}] & D \end{tikzcd}

{\tikzcdset{foo/.style={two heads},every arrow/.append style={thick},every label/.append style={inner sep=2pt}}\begin{tikzcd}[arrows=foo] A \arrow[r,"f"] & B \end{tikzcd}}

\begin{tikzcd}[baseline=(current bounding box.north)] A \arrow[d] \\ B \end{tikzcd}

\begin{tikzcd}[baseline=(current bounding box.south east)] A \arrow[d] \\ B \end{tikzcd}

\begin{tikzcd}[baseline=(Q.south)] |[alias=Q]| A \arrow[d] \\ B \end{tikzcd}

\begin{tikzcd}[baseline=(\tikzcdmatrixname-1-1.north east)] A \arrow[d] \\ B \end{tikzcd}

\begin{tikzcd}[baseline=(L.north)] A \arrow[d,"f"{name=L}] \\ B \end{tikzcd}

\begin{tikzcd}[baseline=(\tikzcdmatrixname.north)] A \arrow[d] \\ B \end{tikzcd}

\begin{tikzcd}[baseline=1ex] A \arrow[r] & B \end{tikzcd}

\begin{tikzcd} |[draw,rounded corners=2pt]| A \arrow[r] & |[sharp corners,draw=blue]| BBB \arrow[r] & |[rounded corners,draw]| \frac{a}{b} \end{tikzcd}

\begin{tikzcd}[line width=1pt,cells={nodes={draw}}] A \arrow[r] & B \end{tikzcd}

\begin{tikzcd}[line width=1pt] |[draw]| A \arrow[r] \arrow[d] & B \\ C \end{tikzcd}

\begin{tikzcd}[column sep=large] A \arrow[r,"a"{anchor=mid}] & B \arrow[r,"b"{anchor=base east}] & C \arrow[r,"c"{anchor=30}] & D \arrow[r,"d"{anchor=real center,marking}] & E \arrow[r,"e"{anchor=text}] & F \end{tikzcd}

\begin{tikzcd} A \arrow[to=\tikzcdmatrixname-1-2] & B \\ C \arrow[from=\tikzcdmatrixname-1-2] \end{tikzcd}

\begin{tikzcd}[sep=large] A \arrow[r,start anchor=mid,end anchor=mid west] \arrow[d,start anchor=base west,end anchor=base east] & B \\ C \arrow[r,start anchor=real center,end anchor=north west] \arrow[ru,start anchor=south east,end anchor=text] & D \end{tikzcd}

\begin{tikzcd}[sep=large] A \arrow[r,start anchor=-30,end anchor=200] \arrow[d,start anchor=10.7,end anchor=90] & B \\ C \end{tikzcd}

{\tikzcdset{row sep/normal=1cm,column sep/large={2cm,between origins}}\begin{tikzcd}[column sep=large] A \arrow[r] \arrow[d] & BBBBBBBB \\ C \end{tikzcd}}

\begin{tikzcd}[every matrix/.append style={inner sep=2pt}] A \arrow[r,"f"] & B \end{tikzcd}

\begin{tikzcd}[column sep=large] A \arrow[r,squiggly,Rightarrow] & B \arrow[r,squiggly,thick,hook] & C \arrow[r,squiggly,dashed,crossing over] & D \arrow[r,bend left,squiggly,crossing over,"f"] & E \end{tikzcd}

\begin{tikzcd} A \arrow[r,shift right] \arrow[r,shift right=.5,shift left=2pt] & B \arrow[r,shift left=1pt,shift left=1pt] & C \arrow[r,shorten <=1pt,shorten >=2pt,crossing over,bend left] & D \end{tikzcd}

\begin{tikzcd}[math mode=false,column sep=large] A \arrow[r,"x"{font=\small},"y"'{near start,text=red}] & B \arrow[r,phantom,"z"{anchor=north}] & C \end{tikzcd}

\begin{tikzcd} A \arrow[r,near start,swap,"a","b"',"c"{midway}] & B \arrow[r,sloped,"d"] & C \arrow[r,pos=.2,"e"] & D \end{tikzcd}
\end{document}
